\documentclass[a4paper, 12pt, twoside]{article}
\usepackage[utf8]{inputenc}
\usepackage[T1]{fontenc}
\usepackage[francais]{babel}
\usepackage{lmodern}
\usepackage{ae,aecompl}
\usepackage[top=2.5cm, bottom=2cm,
            left=3cm, right=2.5cm,
            headheight=15pt]{geometry}
\usepackage{graphicx}
\usepackage{eso-pic}
\usepackage{array}
\usepackage{hyperref}
\usepackage{booktabs}
\usepackage{amsmath}
\usepackage{eurosym}
%%%%%%%%%%%%%%%%%%%%%%%%%%%%%%%%%%%%%%%%
%    Page de garde (Pagedegarde.tex)   %
%%%%%%%%%%%%%%%%%%%%%%%%%%%%%%%%%%%%%%%%
% Dorian Depriester, 2014

\makeatletter
\def\@ecole{école}
\newcommand{\ecole}[1]{
  \def\@ecole{#1}
}

\def\@entreprise{Nom de l'entreprise}
\newcommand{\entreprise}[1]{
  \def\@entreprise{#1}
}

\def\@datedebut{\today}
\newcommand{\datedebut}[1]{
  \def\@datedebut{#1}
}


\def\@datefin{\today}
\newcommand{\datefin}[1]{
  \def\@datefin{#1}
}



\def\@specialite{Spécialité}
\newcommand{\specialite}[1]{
  \def\@specialite{#1}
}

\def\@ED{\'{E}cole Doctorale}
\newcommand{\ED}[1]{
  \def\@ED{#1}
}

\def\@doctorat{Doctorat}
\newcommand{\doctorat}[1]{
  \def\@doctorat{#1}
}

\def\@adresse{Adresse}
\newcommand{\adresse}[1]{
  \def\@adresse{#1}
}

\def\@directeur{directeur}
\newcommand{\directeur}[1]{
  \def\@directeur{#1}
}

\def\@encadrant{encadrant}
\newcommand{\encadrant}[1]{
  \def\@encadrant{#1}
}
\def\@membrea{Membre}
\newcommand{\membrea}[1]{
  \def\@membrea{#1\\}
}
\def\@membreb{Membre}
\newcommand{\membreb}[1]{
  \def\@membreb{#1\\}
}
\def\@membrec{Membre}
\newcommand{\membrec}[1]{
  \def\@membrec{#1\\}
}
\def\@membred{Membre}
\newcommand{\membred}[1]{
  \def\@membred{#1\\}
}
\def\@membree{Membre}
\newcommand{\membree}[1]{
  \def\@membree{#1\\}
}





\def\@juryb{}{}{}
\newcommand{\juryb}[3]{
  \def\@juryb{#1,	& #2	& #3\\}
}
\def\@juryc{}{}{}
\newcommand{\juryc}[3]{
  \def\@juryc{#1,	& #2	& #3\\}
}
\def\@juryd{}{}{}
\newcommand{\juryd}[3]{
  \def\@juryd{#1,	& #2	& #3\\}
}
\def\@jurye{}{}{}
\newcommand{\jurye}[3]{
  \def\@jurye{#1,	& #2	& #3\\}
}
\def\@juryf{}{}{}
\newcommand{\juryf}[3]{
  \def\@juryf{#1,	& #2	& #3\\}
}
\def\@juryg{}{}{}
\newcommand{\juryg}[3]{
  \def\@juryg{#1,	& #2	& #3\\}
}
\def\@juryh{}{}{}
\newcommand{\juryh}[3]{
  \def\@juryh{#1,	& #2	& #3\\}
}
\def\@juryi{}{}{}
\newcommand{\juryi}[3]{
  \def\@juryi{#1,	& #2	& #3\\}
}
\makeatother

\newcommand\BackgroundPic{%
	\put(0,0){%
		\parbox[b][\paperheight]{\paperwidth}{%
			\includegraphics[height=0.45\paperheight]{bordure.png}%
			\vfill
		}
	}
}
\newcommand\EtiquetteThese{%
	\put(0,0){%
		\parbox[t][\paperheight]{\paperwidth}{%
			\hfill
			%\colorbox{blue}{		
				\begin{minipage}[b]{2em}
					\includegraphics[width=4.0\textwidth]{logo_miage.png}\\					
					%\centering\Huge\textcolor{white}{M\\I\\A\\G\\E\\}
					\vspace{0.2cm}
				\end{minipage}
			%}
		}
	}
}

\makeatletter
\newcommand{\pagedegarde}{
\newgeometry{top=2.5cm, bottom=1cm, left=2cm, right=1cm}
\AddToShipoutPicture*{\BackgroundPic}
%\AddToShipoutPicture*{\EtiquetteThese}
  \begin{titlepage}
  \centering
      \includegraphics[width=0.6\textwidth]{logo_Paris_Nanterre_couleur_RVB.png}
      \hfill
      $\ $\\
      %\includegraphics[width=0.20\textwidth]{logo_entreprise.png}\\
    \vspace{1cm}
      {\Large Licence MIASHS première année}\\
    \vspace{1cm}
      {\huge 
      	{\bfseries Rapport de projet informatique}\\
    \vspace{0.5cm}}
      	$\ $\\
    \vspace{1cm}
   		
    \vspace{1cm}
    	{\huge\color[rgb]{0,0,1} \bfseries{\@title}}\\
    \vspace{0.5cm}
    %{\bfseries Entreprise d'accueil : \@entreprise}\\
    {\bfseries Projet réalisé du \@datedebut\ au \@datefin}\\
    %	{\Large{\bfseries Spécialité doctorale ``\@specialite''}}\\
    \vspace{2cm}
    $\ $\\
    \vspace{0.5cm}
    $\ $\\
    \vspace{0.5cm}
    %	le \@date \\
    \vfill
     %  {\LARGE \color[rgb]{0,0,1} \bfseries{\@title}} \\
    %\vfill
      %  Directeur de thèse : {\bfseries \@directeur}\\
       % Co-encadrant de thèse : {\bfseries \@encadrant}\\
    %\vfill
	\begin{tabular}{>{\bfseries}lll}
		\large Membres du groupe\\
		\vspace{0.15cm}\\
		\@membrea
		\@membreb
		\@membrec
		\@membred
		\@membree
		%\@jurye
		%\@juryf
		%\@juryg
		%\@juryh
		%\@juryi
	\end{tabular}
	%\includegraphics[width=0.20\textwidth]{logo_entreprise.png}\
	\vfill
	
	%\@adresse
  \end{titlepage}




\restoregeometry  
}


\title{Quant Terminal : Simulateur d'impact d'événements économiques sur un portefeuille financier}
\entreprise{Université Paris Nanterre --- L1 Double Licence MIASHS \& Économie-Gestion}
\datedebut{3 avril 2026}
\datefin{23 mai 2026}

\membrea{Quentin Geldreich --- 45009207}
\membreb{Lucas Doazan --- 45004502}
\membrec{Evan Saadi --- 45019089}
\membred{Alex Ruimy --- 45015710}

\begin{document}
\pagedegarde

% ============================================================
\section*{Remerciements}
% ============================================================

Nous tenons à remercier notre professeur pour ce cours d'informatique qui nous a permis
de découvrir la programmation en Python et de mener à bien ce projet.

Nous remercions également les équipes derrière les outils que nous avons utilisés :
Streamlit, OpenAI et Yahoo Finance, sans lesquels ce projet n'aurait pas été possible.

Enfin, nous remercions l'intelligence artificielle Claude, développée par Anthropic,
qui a joué un rôle central dans le développement de ce projet. Nous tenons à être
transparents là-dessus : Claude a été bien plus qu'un simple correcteur de syntaxe.
C'est avec lui qu'on a débogué les simulations irréalistes, géré les bugs de sidebar,
et construit l'architecture du projet. Ce rapport parle d'un projet de programmation
assistée par IA, et c'est vraiment ce que c'était.

\newpage
\tableofcontents
\newpage

% ============================================================
\section{Introduction}
% ============================================================

\subsection*{Contexte}

Aujourd'hui, les marchés financiers sont au cœur de l'économie mondiale. Les crises,
les décisions des banques centrales ou les avancées technologiques ont des répercussions
directes sur la valeur des actifs financiers. Pourtant, comprendre et anticiper ces
impacts reste difficile pour quelqu'un qui n'est pas expert en finance.

En tant qu'étudiants en première année de double licence MIASHS et Économie-Gestion
à l'Université Paris Nanterre, nous avons un intérêt commun pour l'investissement et
la finance de marché. Ce projet est né de l'envie de combiner nos deux domaines
d'étude : les mathématiques et l'informatique d'un côté, l'économie et la finance
de l'autre.

\subsection*{Problématique}

Comment permettre à un étudiant ou à un particulier, sans connaissances avancées en
finance, de simuler l'impact d'un événement économique sur un portefeuille
d'investissement de manière réaliste et pédagogique ?

\subsection*{Notre solution : Quant Terminal}

Quant Terminal est une application web interactive développée en Python avec le
framework Streamlit. L'utilisateur décrit un événement économique en langage naturel
(une hausse des taux directeurs, une crise technologique, un choc géopolitique) et
observe l'impact simulé sur un portefeuille financier personnalisé.

L'application s'appuie sur l'intelligence artificielle (OpenAI GPT) pour analyser le
scénario et estimer les chocs sur chaque actif, sur Yahoo Finance pour récupérer les
vrais prix de marché en temps réel, et sur des modèles mathématiques (mouvement
brownien géométrique, Monte-Carlo) pour simuler l'évolution du portefeuille jour
après jour.

\subsection*{Plan du rapport}

Ce rapport présente d'abord notre environnement de travail et les outils utilisés,
puis l'architecture et le fonctionnement du projet avec les modèles mathématiques
sous-jacents. Nous détaillons ensuite les fonctionnalités réalisées, les difficultés
rencontrées, et nous terminons par un bilan de cette expérience.

\newpage

% ============================================================
\section{Environnement de travail}
% ============================================================

\subsection{Outils de développement}

Nous avons développé le projet en Python 3.10 sous Visual Studio Code. La gestion
des dépendances se fait via un fichier \texttt{requirements.txt}, et la clé API
OpenAI est stockée dans un fichier \texttt{.env} pour des raisons de sécurité.
Ce fichier n'est jamais envoyé sur GitHub.

L'outil qui a le plus accompagné notre développement est \textbf{Claude Code}
(Anthropic), un assistant IA intégré directement dans Visual Studio Code. Nous avons
ouvert notre dossier de projet entier dans Claude Code, ce qui lui permettait d'avoir
une vision globale du code. En pratique, on décrivait une fonctionnalité en français,
Claude proposait un code, on testait, ça ne marchait pas toujours, on corrigeait
ensemble. C'est ce cycle qui a structuré tout notre développement.

\subsection{Collaboration et versionnage}

Nous avons utilisé \textbf{Git} et \textbf{GitHub} pour gérer les versions du projet.
Le dépôt est public et accessible à l'adresse suivante :

\begin{center}
\url{https://github.com/00quentin/projet-info-ia-L1-QG-LD-AR-ES}
\end{center}

Les commits étaient principalement réalisés par les membres chargés du développement.
Les décisions importantes étaient discutées collectivement par les quatre membres.

\subsection{Déploiement}

L'application est déployée en ligne via \textbf{Streamlit Cloud}, accessible
gratuitement depuis n'importe quel navigateur. Un problème de mise en veille
automatique après quelques minutes d'inactivité a été résolu grâce à
\textbf{Uptime Robot}, un service qui envoie une requête au site toutes les cinq
minutes pour le maintenir actif.

\subsection{Organisation du groupe}

Le groupe est composé de quatre membres. La collaboration s'est organisée en deux
phases : plusieurs séances en présentiel durant les TD pour définir les objectifs et
répartir le travail, puis une longue phase de développement à distance, avec une
réunion finale pour tout mettre en forme.

\newpage

% ============================================================
\section{Description du projet et objectifs}
% ============================================================

\subsection{Cahier des charges initial}

L'idée de départ était simple : permettre à un utilisateur de saisir un événement
économique et d'observer son impact sur un portefeuille. En cours de développement,
nous avons progressivement enrichi le projet avec de nombreuses fonctionnalités
supplémentaires qui n'étaient pas prévues au départ.

Les fonctionnalités prévues dès le début étaient :

\begin{itemize}
    \item Saisie d'un événement économique en langage naturel
    \item Analyse de l'événement par intelligence artificielle
    \item Simulation de l'impact sur un portefeuille personnalisable
    \item Visualisation des résultats sous forme de graphiques interactifs
    \item Affichage de métriques financières (rendement, risque)
\end{itemize}

\subsection{Architecture du projet}

Le projet est organisé en modules distincts pour séparer clairement la logique métier
de l'interface utilisateur. La règle principale que nous avons suivie : le dossier
\texttt{core/} ne doit jamais importer Streamlit. Cela permet de tester la logique
indépendamment de l'interface, et d'éviter que tout le projet s'effondre quand
Streamlit change de version.

\begin{table}[h]
\begin{center}
\begin{tabular}{|l|p{9cm}|}
\hline
\textbf{Fichier / Dossier} & \textbf{Rôle} \\
\hline
\texttt{app.py} & Point d'entrée principal, orchestre les onglets et la navigation \\
\hline
\texttt{core/} & Logique métier pure (calculs, métriques, historique) sans import Streamlit \\
\hline
\texttt{components/} & Composants d'interface réutilisables (sidebar, graphiques, notifications) \\
\hline
\texttt{views/} & Un fichier par onglet (dashboard, portefeuille, comparaison, etc.) \\
\hline
\texttt{handlers.py} & Orchestration des lancements de simulation et de backtest \\
\hline
\texttt{ia\_bot.py} & Appels à l'API OpenAI avec validation des réponses \\
\hline
\texttt{market\_data.py} & Récupération des prix via Yahoo Finance avec mise en cache \\
\hline
\texttt{simulation.py} & Moteur de simulation (probabiliste, historique, Monte-Carlo) \\
\hline
\texttt{schemas.py} & Validation des réponses de l'IA via Pydantic \\
\hline
\texttt{config.py} & Toutes les constantes du projet (actifs, profils, bornes) \\
\hline
\end{tabular}
\end{center}
\caption{Architecture modulaire de Quant Terminal}
\end{table}

\subsection{Fonctionnement général}

Voici les étapes qu'un utilisateur suit pour réaliser une simulation :

\begin{enumerate}
    \item Choisir le mode : \textit{Simulation prospective} (scénario fictif) ou
    \textit{Backtest historique} (crise réelle passée)
    \item Décrire un événement économique en langage naturel dans la sidebar
    \item Sélectionner les actifs à inclure (jusqu'à 23 actifs disponibles)
    \item Définir le capital de départ et le profil d'investisseur
    \item Cliquer sur \textit{Lancer la simulation}
    \item L'IA analyse le scénario et estime les chocs sur chaque actif,
    validés par Pydantic pour éviter les valeurs aberrantes
    \item Le moteur calcule l'évolution du portefeuille jour après jour
    \item Les résultats s'affichent : graphiques, métriques, analyse IA, alertes
\end{enumerate}

\subsection{Le mouvement brownien géométrique}

Le moteur de simulation repose sur le \textbf{mouvement brownien géométrique} (GBM),
un modèle standard en finance quantitative, utilisé notamment dans la formule de
Black-Scholes pour la valorisation des options. L'équation différentielle stochastique
qui le décrit est :

\begin{equation}
dS_t = \mu \, S_t \, dt + \sigma \, S_t \, dW_t
\end{equation}

où $S_t$ est la valeur de l'actif au temps $t$, $\mu$ le drift (tendance directionnelle
imposée par le choc estimé par l'IA), $\sigma$ la volatilité historique annualisée
récupérée via Yahoo Finance sur 12 mois, et $dW_t$ un processus de Wiener représentant
le bruit aléatoire du marché.

En pratique, à chaque jour de cotation simulé, la variation de chaque actif est
calculée via la solution exacte du GBM :

\begin{equation}
S_{t+1} = S_t \cdot \exp\!\left[\left(\mu - \frac{\sigma^2}{2}\right)\Delta t
+ \sigma \sqrt{\Delta t} \cdot \varepsilon\right]
\end{equation}

où $\varepsilon \sim \mathcal{N}(0,1)$ est un nombre aléatoire tiré d'une loi normale
standard (généré via NumPy). Le terme $-\sigma^2/2$ est une correction d'Itô qui
compense le biais introduit par la nature multiplicative du processus.

Nous avons choisi ce modèle pour trois raisons : il garantit que les prix restent
toujours positifs (contrairement à un modèle additif), il est calibrable directement
à partir des données Yahoo Finance, et c'est le modèle de référence académique en
finance stochastique.

\textbf{Mode Monte-Carlo :} en activant ce mode, le processus est répété 50 fois avec
des tirages $\varepsilon$ différents. On obtient ainsi une enveloppe de trajectoires
possibles plutôt qu'une courbe unique. Les zones colorées sur les graphiques
représentent les percentiles 5\% et 95\% de ces 50 simulations.

\newpage

% ============================================================
\section{Bibliothèques, Outils et technologies}
% ============================================================

\begin{table}[h]
\begin{center}
\begin{tabular}{|l|p{5.5cm}|p{4cm}|}
\hline
\textbf{Bibliothèque} & \textbf{Rôle dans le projet} & \textbf{Pourquoi ce choix} \\
\hline
Streamlit & Framework principal pour l'interface web & Simple en Python pur, déploiement gratuit \\
\hline
OpenAI API & Analyse des scénarios par IA & GPT comprend le langage naturel en contexte financier \\
\hline
yfinance & Prix réels via Yahoo Finance & Gratuit, fiable, supporte tous nos actifs \\
\hline
Plotly & Graphiques interactifs & Rendu professionnel, natif dans Streamlit \\
\hline
Pydantic & Validation des réponses de l'IA & Évite les valeurs aberrantes type Or +1 000 000\% \\
\hline
NumPy & Calculs du GBM et Monte-Carlo & Référence pour les maths scientifiques en Python \\
\hline
Pandas & Séries temporelles journalières & Facilite la manipulation des données de marché \\
\hline
ReportLab & Génération du rapport PDF & Création de PDF directement en Python \\
\hline
python-dotenv & Gestion de la clé API & Sécurise la clé OpenAI hors du code source \\
\hline
\end{tabular}
\end{center}
\caption{Bibliothèques utilisées dans Quant Terminal}
\end{table}

\subsection{Note sur Streamlit}

Streamlit est une bibliothèque Python qui transforme automatiquement chaque widget
Python (bouton, slider, graphique) en composant web interactif, sans avoir besoin
d'écrire du HTML ou du JavaScript. C'est ce qui nous a permis, avec notre niveau de
L1, de créer une vraie application web.

Sa limite principale : Streamlit recrée entièrement la page à chaque interaction
utilisateur. Ça simplifie beaucoup de choses, mais ça rend certains comportements
d'interface très difficiles à contrôler, comme on l'a découvert avec le bouton
de la sidebar.

\subsection{Note sur la validation Pydantic}

Pydantic est une bibliothèque de validation de données en Python. Dans notre projet,
elle sert à vérifier que les chocs estimés par l'IA sont dans des plages réalistes
avant d'être injectés dans le GBM. Par exemple, le schéma de validation pour le
S\&P 500 impose une borne entre -55\% et +45\%, calibrée sur les pires crises
historiques (2008 : -57\% sur 18 mois). Si l'IA dépasse ces bornes, Pydantic lève
une erreur, la valeur est corrigée et un avertissement est affiché.

\newpage

% ============================================================
\section{Travail réalisé}
% ============================================================

\subsection{Fonctionnalités réalisées}

\begin{itemize}
    \item[$\checkmark$] \textbf{Simulation prospective} : saisie d'un événement fictif,
    analyse IA, simulation GBM sur 1 à 3 scénarios comparables simultanément
    \item[$\checkmark$] \textbf{Backtest historique} : rejouer de vraies crises
    (COVID-19, 2008, Choc pétrolier 1973, etc.) avec les données réelles Yahoo Finance
    \item[$\checkmark$] \textbf{Mode comparaison} : graphique d'évolution comparée et
    tableau de métriques côte à côte pour 2 ou 3 scénarios
    \item[$\checkmark$] \textbf{Mode Monte-Carlo} : 50 simulations GBM parallèles avec
    enveloppe probabiliste (percentiles 5-95\%)
    \item[$\checkmark$] \textbf{Calibration historique IA} : l'IA s'inspire des
    amplitudes réelles de crises passées pour calibrer ses estimations de chocs
    \item[$\checkmark$] \textbf{Portefeuille personnalisable} : capital, profil
    d'investisseur (Prudent / Équilibré / Agressif / Personnalisé), allocation manuelle
    \item[$\checkmark$] \textbf{Actifs personnalisés} : ajout de n'importe quel actif
    via son ticker Yahoo Finance (TSLA, BTC-USD, etc.)
    \item[$\checkmark$] \textbf{Prix en temps réel} : vrais prix Yahoo Finance avec
    cache d'une heure et bouton de rafraîchissement manuel
    \item[$\checkmark$] \textbf{Métriques institutionnelles} : Volatilité annualisée,
    Sharpe Ratio, Max Drawdown, VaR 95\% avec alertes automatiques
    \item[$\checkmark$] \textbf{Historique des simulations} : sauvegarde automatique
    des 10 dernières simulations en JSON, persistante entre les sessions
    \item[$\checkmark$] \textbf{Export CSV} : données journalières brutes par actif
    \item[$\checkmark$] \textbf{Export PDF} : rapport complet rédigé par l'IA avec
    analyse, métriques et recommandations
    \item[$\checkmark$] \textbf{Chatbot IA} : analyste financier conversationnel
    accessible via une bulle flottante
    \item[$\checkmark$] \textbf{Académie pédagogique} : 7 modules sur la finance de
    marché (lecture des données, macro-économie, cas historiques, stratégies, etc.)
    \item[$\checkmark$] \textbf{Mode sombre / clair} : interface adaptable
    \item[$\checkmark$] \textbf{Partage par URL} : paramètres de simulation encodés
    dans l'URL pour partager une configuration
    \item[$\checkmark$] \textbf{Ticker en direct} : S\&P 500, VIX, Or, Pétrole,
    BTC, EUR/USD mis à jour en temps réel
\end{itemize}

\subsection{Fonctionnalités non réalisées}

Certaines idées ont été abandonnées en cours de route. Nous voulions une interface
visuellement plus proche des plateformes fintech professionnelles, mais les contraintes
de Streamlit nous en ont empêchés. Chaque effet visuel avancé nécessitait des
contournements CSS complexes qu'on n'avait pas anticipés.

Des fonctionnalités plus ambitieuses comme l'authentification utilisateur ou une base
de données pour l'historique ont été écartées : elles dépassaient le cadre d'un projet
de première année et le temps disponible.

\subsection{Répartition réelle du travail}

\begin{table}[h]
\begin{center}
\begin{tabular}{|l|p{10cm}|}
\hline
\textbf{Membre} & \textbf{Tâches réalisées} \\
\hline
Quentin Geldreich & Développement principal : architecture générale, simulation GBM,
intégration OpenAI et Yahoo Finance, interface Streamlit, déploiement en ligne \\
\hline
Lucas Doazan & Co-développement, tests réguliers, suggestions de fonctionnalités,
participation aux décisions techniques et à la revue de code \\
\hline
Evan Saadi & Rédaction du rapport, recherches sur les concepts financiers (GBM,
métriques de risque), tests fonctionnels de l'application, retours utilisateur \\
\hline
Alex Ruimy & Rédaction du rapport, tests de l'application sur différents scénarios,
retours sur l'expérience utilisateur, suggestions d'amélioration de l'interface \\
\hline
\end{tabular}
\end{center}
\caption{Répartition réelle du travail au sein du groupe}
\end{table}

Toutes les décisions importantes (fonctionnalités à ajouter, scénarios à tester,
orientation du projet) étaient discutées collectivement. Les membres chargés du
rapport ont aussi joué un rôle de testeurs réguliers, en signalant les bugs et en
vérifiant que les résultats de simulation semblaient cohérents.

\newpage

% ============================================================
\section{Difficultés rencontrées}
% ============================================================

\subsection{Les limitations de Streamlit}

La principale difficulté technique du projet a été de travailler avec Streamlit sur
des interfaces un peu avancées. Le framework est excellent pour démarrer rapidement,
mais il montre vite ses limites quand on veut personnaliser le comportement de
certains composants.

Le cas le plus frustrant a été le bouton d'ouverture et de fermeture de la sidebar.
Ce bouton fonctionnait parfaitement au départ, puis disparaissait définitivement
après le premier clic de fermeture, rendant la sidebar totalement inaccessible.
On a mis plusieurs jours à comprendre pourquoi : Streamlit réécrit entièrement le
DOM (la structure de la page web) en JavaScript à chaque interaction, écrasant
systématiquement nos modifications CSS. On a essayé plusieurs approches : injecter
du JavaScript pour intercepter le clic, cibler le bouton par différents sélecteurs
CSS, modifier le comportement via les query params. Rien ne tenait sur la durée.
La solution finale a été de forcer la sidebar à rester toujours visible via du CSS
\texttt{!important} et de cacher le bouton de fermeture. Ce n'est pas élégant, mais
c'est la seule approche stable qu'on ait trouvée.

En fait, chaque effet visuel un peu avancé nécessitait des contournements qu'on
n'avait pas anticipés au départ. Le design qu'on imaginait au début aurait
probablement nécessité de travailler avec React ou une autre technologie.

\subsection{La fiabilité des simulations financières}

Lors des premières simulations, l'IA renvoyait des valeurs complètement irréalistes.
L'exemple le plus marquant : dans l'une de nos premières versions, l'Or affichait
une variation de \textbf{+1 222 891\%} lors d'un scénario de crise. Ce genre de
résultat rendait l'application absurde et inutilisable.

Comprendre pourquoi ça arrivait a pris du temps. Le problème venait du fait qu'on
demandait à l'IA d'estimer des pourcentages de variation sans lui donner de cadre de
référence. Sans contrainte, GPT optimisait pour la "dramatisation" du scénario plutôt
que pour le réalisme. On a résolu ça en deux étapes : d'abord en précisant dans le
prompt les ordres de grandeur attendus, puis en ajoutant une validation Pydantic avec
des bornes calibrées sur les vraies crises historiques. Depuis, les simulations
produisent des résultats cohérents.

\subsection{Le déploiement sur Streamlit Cloud}

Le déploiement a posé plusieurs problèmes qu'on n'avait pas anticipés. Le principal :
l'application se mettait en veille automatiquement après quelques minutes d'inactivité,
obligeant les visiteurs à attendre 30 secondes à chaque rechargement. On a contourné
ça avec Uptime Robot, un service gratuit qui envoie une requête au site toutes les
cinq minutes.

La configuration de la clé API OpenAI a aussi nécessité une manipulation spécifique :
sur Streamlit Cloud, le fichier \texttt{.env} local n'existe pas. Il faut renseigner
les secrets directement dans l'interface de déploiement.

\subsection{Le comportement du chatbot}

Au départ, le chatbot répondait à n'importe quelle question sans se limiter à la
finance. On a dû modifier le prompt système pour le recadrer sur son domaine. Ce
travail de "prompt engineering" (la manière dont on formule les instructions données
à un modèle de langage) a été plus complexe qu'attendu : trop restrictif, le chatbot
refusait des questions légitimes ; trop permissif, il dérivait vers des sujets sans
rapport. On a fini par trouver un équilibre qui fonctionne bien.

\newpage

% ============================================================
\section{Bilan}
% ============================================================

\subsection{Conclusion}

Honnêtement, on ne pensait pas arriver à un résultat comme celui-là quand on a
commencé. L'idée de départ était simple : saisir un événement, voir l'impact sur un
portefeuille. Un mois et demi plus tard on a une application déployée en ligne avec
17 fonctionnalités, des graphiques interactifs, un modèle Monte-Carlo et un chatbot
IA intégré. C'est difficile de ne pas en être fiers.

On a appris des tas de choses sans vraiment s'en rendre compte. Au détour d'un bug
à corriger on comprenait comment fonctionne le cache en Python ; en essayant de rendre
les simulations réalistes on s'est retrouvés à lire des articles sur les crises de
1973 et sur le modèle de Black-Scholes ; en déployant sur Streamlit Cloud on a
découvert ce qu'est une variable d'environnement et pourquoi on ne commit jamais une
clé API sur GitHub. C'est ce genre d'apprentissage concret qu'on n'aurait pas eu en
restant sur des exercices académiques.

Sur l'IA : Claude a été un vrai partenaire de développement, pas juste un
autocorrecteur. Ce n'est pas "l'IA a fait le projet à notre place". C'est plus proche
de : l'IA nous a permis d'aller bien au-delà de notre niveau technique actuel. Sans
elle, on aurait probablement produit un script Python basique. Avec elle, on a produit
une vraie application. Mais sans notre direction, nos tests et nos corrections, et il
y en a eu beaucoup, Claude aurait généré des simulations avec l'or à +1 222 891\% et
une sidebar qui disparaît au premier clic.

Si on devait refaire ce projet, on passerait davantage de temps à planifier
l'architecture avant de commencer à coder. On a beaucoup ajouté de fonctionnalités
au fil de l'eau, ce qui a parfois créé des incohérences qu'on a dû corriger ensuite.
On choisirait aussi probablement une technologie plus flexible que Streamlit pour le
design.

\subsection{Perspectives}

On a envie de continuer ce projet après le cours. Ce qui manque le plus au modèle
actuel, c'est la corrélation entre actifs : aujourd'hui chaque actif est simulé
indépendamment, alors que dans la réalité l'or et le dollar bougent souvent en sens
inverse, tout comme les obligations par rapport aux actions. Ce serait la prochaine
étape logique côté maths. On aimerait aussi remplacer la sauvegarde en JSON par une
vraie base de données, avec un système d'authentification pour que chaque utilisateur
garde son propre historique. Et si on recommençait depuis zéro, on choisirait
probablement une technologie plus flexible que Streamlit pour l'interface.

\newpage

% ============================================================
\section{Bibliographie}
% ============================================================

\renewcommand{\bibname}{}
\renewcommand{\refname}{}
\begin{thebibliography}{9}
   \bibitem[BLACK73]{bs} F. Black, M. Scholes,
   {\it The Pricing of Options and Corporate Liabilities},
   Journal of Political Economy, Vol. 81, No. 3, 1973
   \bibitem[HULL]{hull} J. Hull,
   {\it Options, Futures, and Other Derivatives}, 11e édition, Pearson, 2022
   \bibitem[ITO]{ito} K. Itô,
   {\it On Stochastic Differential Equations},
   Memoirs of the American Mathematical Society, 1951
\end{thebibliography}

\newpage

% ============================================================
\section{Webographie}
% ============================================================

\begin{thebibliography}{9}
   \bibitem[ST]{st} Documentation officielle Streamlit,
   \url{https://docs.streamlit.io}
   \bibitem[OAI]{oai} Documentation OpenAI API,
   \url{https://platform.openai.com/docs}
   \bibitem[YF]{yf} yfinance sur PyPI,
   \url{https://pypi.org/project/yfinance}
   \bibitem[PD]{pd} Documentation Pydantic v2,
   \url{https://docs.pydantic.dev}
   \bibitem[PL]{pl} Documentation Plotly Python,
   \url{https://plotly.com/python}
   \bibitem[GBM]{gbm} Wikipédia, Mouvement brownien géométrique,
   \url{https://fr.wikipedia.org/wiki/Mouvement_brownien_g\%C3\%A9om\%C3\%A9trique}
   \bibitem[GH]{gh} Dépôt GitHub du projet,
   \url{https://github.com/00quentin/projet-info-ia-L1-QG-LD-AR-ES}
   \bibitem[UR]{ur} Uptime Robot, monitoring de disponibilité,
   \url{https://uptimerobot.com}
\end{thebibliography}

\newpage

% ============================================================
\appendix
\makeatletter
\def\@seccntformat#1{Annexe~\csname the#1\endcsname:\quad}
\makeatother

% ---- ANNEXE A ----
\section{Cahier des charges}

Le cahier des charges initial comprenait les fonctionnalités suivantes :

\begin{itemize}
    \item Saisie d'un événement économique en langage naturel
    \item Analyse de l'événement par intelligence artificielle (OpenAI GPT)
    \item Sélection des actifs et du profil d'investisseur
    \item Simulation de l'évolution d'un portefeuille sur un horizon défini
    \item Affichage de graphiques interactifs et de métriques financières
\end{itemize}

Les fonctionnalités suivantes ont été ajoutées hors cahier des charges initial, au fil
du développement : mode backtest historique, comparaison de scénarios, Monte-Carlo,
calibration historique IA, export CSV/PDF, chatbot IA, académie pédagogique, mode
sombre, partage par URL, ticker en direct, actifs personnalisés via ticker Yahoo.

% ---- ANNEXE B ----
\section{Exemple d'exécution du projet}

Cette annexe présente une session complète d'utilisation de Quant Terminal.
Nous avons simulé deux scénarios économiques opposés sur le même portefeuille
de 15 000\euro{} avec un profil Agressif, 20 actifs sélectionnés, un horizon de
130 jours de cotation et le mode Monte-Carlo activé (50 simulations parallèles).

\textbf{Scénario A (Stagflation européenne) :} la BCE monte ses taux à 5,5\%
malgré une récession, l'euro chute, les marchés actions européens s'effondrent
mais l'or et le pétrole résistent en valeur refuge.

\textbf{Scénario B (Boom IA et semi-conducteurs asiatiques) :} la Chine lève ses
restrictions sur les semi-conducteurs, les valeurs tech explosent, les actifs
risqués surperforment massivement.

Résultats : Scénario A perd \textbf{-10,81\%} (13 379\euro{}),
Scénario B gagne \textbf{+24,32\%} (18 648\euro{}). L'écart de 35,12 points entre
les deux trajectoires illustre la puissance de la fonctionnalité de comparaison.

\bigskip
\noindent\textbf{Étape 1 : Page d'accueil et configuration}

\begin{figure}[h]
\centering
\includegraphics[width=0.85\textwidth]{fig_accueil1.png}
\caption{Page d'accueil en mode sombre. Le ticker en haut affiche les prix en direct
(S\&P 500, VIX, Or, BTC) et un guide en 4 étapes explique comment démarrer}
\end{figure}

\begin{figure}[h]
\centering
\includegraphics[width=0.85\textwidth]{fig_accueil2.png}
\caption{Dashboard vide avant la première simulation, avec presets de scénarios
pour aider l'utilisateur à démarrer}
\end{figure}

\begin{figure}[h]
\centering
\includegraphics[width=0.85\textwidth]{fig_accueil3.png}
\caption{Scénarios rapides prêts à l'emploi. Un clic charge la configuration complète
sans avoir à remplir le formulaire manuellement}
\end{figure}

\begin{figure}[h]
\centering
\includegraphics[width=0.45\textwidth]{fig_sidebar1.png}
\caption{Sidebar de configuration : mode Simulation prospective, 2 scénarios à
comparer, saisie des événements A et B. Le conseil en jaune incite à détailler le
scénario pour améliorer la qualité de l'analyse}
\end{figure}

\begin{figure}[h]
\centering
\includegraphics[width=0.45\textwidth]{fig_sidebar2.png}
\caption{Options avancées : Monte-Carlo activé (50 simulations), horizon de 130 jours,
prix réels Yahoo Finance, calibration historique IA activée}
\end{figure}

\begin{figure}[h]
\centering
\includegraphics[width=0.45\textwidth]{fig_sidebar3.png}
\caption{Sélection des actifs (20 sur 23) et paramètres du portefeuille : capital de
15 000\euro{}, profil Agressif}
\end{figure}

\bigskip
\noindent\textbf{Étape 2 : Résultats du Dashboard}

\begin{figure}[h]
\centering
\includegraphics[width=0.85\textwidth]{fig_dashboard_1.png}
\caption{Dashboard après simulation du Scénario A : valeur finale 13 379\euro{}
(-10,81\%), calibré sur le Choc pétrolier de 1973}
\end{figure}

\begin{figure}[h]
\centering
\includegraphics[width=0.85\textwidth]{fig_analyseIA.png}
\caption{Analyse macroéconomique rédigée par l'IA. GPT explique les mécanismes de
transmission du scénario avec les estimations d'impact (inflation +8,52\%, taux
directeurs +5,53\%)}
\end{figure}

\begin{figure}[h]
\centering
\includegraphics[width=0.85\textwidth]{fig_evo_vs_marche.png}
\caption{Évolution du portefeuille sur 130 jours comparée au S\&P 500 (-16,18\%)
et au MSCI World (-16,44\%). Notre portefeuille (-10,81\%) surperforme les indices
grâce à la présence d'or et de pétrole}
\end{figure}

\begin{figure}[h]
\centering
\includegraphics[width=0.85\textwidth]{fig_metriques_portefeuille.png}
\caption{Métriques de risque : Volatilité 9,4\% (FAIBLE), Sharpe -2,54 (MAUVAIS),
Max Drawdown -11,6\% (MODÉRÉ), VaR 95\% -1,0\%/jour. L'alerte rouge signale
le rendement/risque défavorable}
\end{figure}

\begin{figure}[h]
\centering
\includegraphics[width=0.85\textwidth]{fig_comparaison_cali.png}
\caption{Comparaison calibration historique vs projection IA libre par actif. Un écart
supérieur à 15 points indique un scénario peu couvert par les données historiques}
\end{figure}

\begin{figure}[h]
\centering
\includegraphics[width=0.85\textwidth]{fig_graphiques.png}
\caption{Simulations Monte-Carlo par catégorie (50 trajectoires). Les zones colorées
montrent l'enveloppe des scénarios possibles. Les actions ont une dispersion bien plus
large que les obligations}
\end{figure}

\begin{figure}[h]
\centering
\includegraphics[width=0.85\textwidth]{fig_graphique_grand.png}
\caption{Monte-Carlo cryptomonnaies : Bitcoin, Ethereum, XRP et Solana sur 130 jours.
La ligne pointillée blanche est la médiane des 50 simulations}
\end{figure}

\begin{figure}[h]
\centering
\includegraphics[width=0.85\textwidth]{fig_performance_actifs.png}
\caption{Performance par actif : Or +15\%, Pétrole +12\%, Défense +7\% en vert,
contre CAC 40 -16\%, S\&P 500 -16\%, NASDAQ -20\% en rouge}
\end{figure}

\bigskip
\noindent\textbf{Étape 3 : Onglet Portefeuille}

\begin{figure}[h]
\centering
\includegraphics[width=0.85\textwidth]{fig_portefeuille_1.png}
\caption{Détail par actif : chaque carte affiche la valeur finale et la performance.
Le contraste est net entre Pétrole (+11,85\%) et NASDAQ (-20,33\%)}
\end{figure}

\begin{figure}[h]
\centering
\includegraphics[width=0.85\textwidth]{fig_portefeuille_2.png}
\caption{Allocation (donut) et bilan net : S\&P 500 à 30\%, Bitcoin à 20\%, NASDAQ
à 15\%. Résultat final : 13 379\euro{} soit -1 620\euro{} sur les 15 000\euro{}
initiaux}
\end{figure}

\bigskip
\noindent\textbf{Étape 4 : Comparaison des deux scénarios}

\begin{figure}[h]
\centering
\includegraphics[width=0.85\textwidth]{fig_comparatifs_portefeuilles.png}
\caption{Évolution comparée : Scénario B (vert, +24,32\%) contre Scénario A (bleu,
-10,81\%). Les trajectoires divergent dès les premiers jours}
\end{figure}

\begin{figure}[h]
\centering
\includegraphics[width=0.85\textwidth]{fig_comparatif_metriques.png}
\caption{Métriques A vs B : Sharpe -2,54 contre +3,99, Max Drawdown -11,60\% contre
-3,50\%. Le scénario B domine sur tous les indicateurs}
\end{figure}

\bigskip
\noindent\textbf{Étape 5 : Export, Historique et fonctionnalités annexes}

\begin{figure}[h]
\centering
\includegraphics[width=0.85\textwidth]{fig_export_pdf.png}
\caption{Section export : données brutes en CSV, rapport PDF généré par l'IA, bouton
de partage qui copie un lien URL avec tous les paramètres de simulation}
\end{figure}

\begin{figure}[h]
\centering
\includegraphics[width=0.85\textwidth]{fig_historique.png}
\caption{Onglet Historique : les deux simulations sont sauvegardées avec date, profil
et performance. L'historique persiste entre les sessions via un fichier JSON}
\end{figure}

\begin{figure}[h]
\centering
\includegraphics[width=0.85\textwidth]{fig_academie_1.png}
\caption{Académie, module Stratégies d'investissement. Le portefeuille 60/40
(actions/obligations) est présenté comme référence institutionnelle depuis les
années 1950}
\end{figure}

\begin{figure}[h]
\centering
\includegraphics[width=0.85\textwidth]{fig_academie_2.png}
\caption{Exemple d'allocation équilibrée : 40\% actions, 25\% obligations, 15\% or,
10\% crypto, 10\% cash}
\end{figure}

\begin{figure}[h]
\centering
\includegraphics[width=0.85\textwidth]{fig_academie_3.png}
\caption{Cas historique de la Crise de 2008 : contexte subprimes, déclenchement,
chute de -57\% sur le S\&P 500 et enseignements}
\end{figure}

\begin{figure}[h]
\centering
\includegraphics[width=0.85\textwidth]{fig_anlyste_ia.png}
\caption{Chatbot IA répondant à \og{}Pourquoi l'Or monte-t-il en période de crise ?\fg{}
Les questions fréquentes à gauche permettent de démarrer sans connaissances préalables}
\end{figure}

\begin{figure}[h]
\centering
\includegraphics[width=0.85\textwidth]{fig_apropos_1.png}
\caption{Page À propos : présentation du projet et contexte académique (L1 MIASHS,
double licence Économie-Gestion, Université Paris Nanterre)}
\end{figure}

\begin{figure}[h]
\centering
\includegraphics[width=0.85\textwidth]{fig_apropos_2.png}
\caption{Stack technique en trois piliers : Frontend/Backend, IA/Données, Hébergement}
\end{figure}

\begin{figure}[h]
\centering
\includegraphics[width=0.85\textwidth]{fig_apropos_3.png}
\caption{Liste des fonctionnalités et section \og{}Sources et limites assumées\fg{} :
Quant Terminal est un outil pédagogique utilisant un GBM avec volatilités supposées
constantes et actifs indépendants}
\end{figure}

\begin{figure}[h]
\centering
\includegraphics[width=0.85\textwidth]{fig_donnees_brut.png}
\caption{Export CSV dans Excel : une colonne par actif, une ligne par jour simulé}
\end{figure}

\begin{figure}[h]
\centering
\includegraphics[width=0.85\textwidth]{fig_pdf.png}
\caption{Rapport PDF généré (4 pages) : scénario, paramètres, synthèse macro de l'IA,
impacts par actif et métriques de risque}
\end{figure}

% ---- ANNEXE C ----
\section{Manuel utilisateur}

\subsection*{Prérequis pour lancer en local}

\begin{itemize}
    \item Python 3.10 ou supérieur installé sur la machine
    \item Un fichier \texttt{.env} à la racine du projet contenant :
    \texttt{OPENAI\_API\_KEY=sk-...}
    \item Dépendances installées via : \texttt{pip install -r requirements.txt}
    \item Lancer l'application avec : \texttt{streamlit run app.py}
\end{itemize}

L'application est accessible directement en ligne (sans installation) à :
\url{https://quant-terminal.streamlit.app}

\subsection*{Guide des onglets}

\begin{itemize}
    \item \textbf{Dashboard} : résultats principaux, valeur finale, analyse IA,
    graphiques d'évolution, Monte-Carlo par catégorie, métriques de risque, alertes
    \item \textbf{Portefeuille} : détail par actif avec performance individuelle et
    donut d'allocation
    \item \textbf{Comparaison} : disponible si 2 ou 3 scénarios sont lancés,
    graphique d'évolution comparée et tableau de métriques côte à côte
    \item \textbf{Historique} : 10 dernières simulations sauvegardées automatiquement,
    persistantes entre les sessions
    \item \textbf{Académie} : 7 modules pédagogiques sur la finance de marché,
    conçus pour être lus dans l'ordre
    \item \textbf{À propos} : présentation du projet, de l'équipe et de la stack
    technique
    \item \textbf{Chatbot IA} : accessible via la bulle violette en bas à droite,
    pour poser des questions financières en langage naturel
\end{itemize}

\end{document}
